% Generated draft: requires mathematical and source review.
\documentclass[11pt]{article}
\usepackage[margin=25mm]{geometry}
\usepackage{fontspec}
\setmainfont{Latin Modern Roman}
\usepackage{amsmath,amssymb,mathtools,unicode-math}
\setmathfont{Latin Modern Math}
\usepackage{enumitem,needspace,xurl,hyperref}
\setlength{\parindent}{0pt}
\setlength{\parskip}{5pt}
\setlength{\emergencystretch}{3em}
\newcommand{\N}{\mathbb{N}}
\newcommand{\Z}{\mathbb{Z}}
\newcommand{\Q}{\mathbb{Q}}
\newcommand{\R}{\mathbb{R}}
\newcommand{\C}{\mathbb{C}}
\newcommand{\F}{\mathbb{F}}
\newcommand{\A}{\mathbb{A}}
\newcommand{\PP}{\mathbb{P}}
\newcommand{\E}{\mathbb{E}}
\newcommand{\eps}{\varepsilon}
\newcommand{\dd}{\,\mathrm{d}}
\begin{document}
\section*{UnsolvedMath ID: 30006200}
\textbf{Source problem code:} catalogue entry
\subsection{Definitions and mathematical statement}
The supplied record for UnsolvedMath ID \(30006200\), source code \texttt{catalogue entry}, is titled ``Optimization Workshop Contents Fragment.'' Its statement consists of a damaged table-of-contents passage beginning with ``Conjecture'' and the page number 416, followed by several talk titles, contributor names, and page numbers. No mathematical assertion accompanies the word ``Conjecture,'' and its association with any listed talk is unspecified.

Consequently, no mathematical definitions or proposition can be recovered from the supplied passage: the intended objects, domains, hypotheses, boundary conditions, quantifiers, and conclusion are absent. The listed talk titles do not determine these missing data. A precise mathematical statement would require an identifiable source passage containing the intended problem.

The supplied statement status is \texttt{unrecoverable}, and the clean-statement field is empty. The registry nevertheless labels the record ``open.'' Its supplied literature assessment explains that this label retains an extraction whose truth cannot be tested because no proposition is identifiable; it does not establish that a particular mathematical conjecture is open. The supplied statement-verification note reports that the official report confirms a damaged table-of-contents extraction and that no separate clean statement was promoted. These are supplied curation claims, not independent verification performed for this draft.
\subsection{Short English statement}
This record preserves a workshop contents fragment, not an identifiable mathematical problem. Its intended conjecture cannot be stated precisely from the supplied evidence. The ``open'' label records supplied curation; the statement itself is marked unrecoverable.
\subsection{Sources}
\begin{itemize}
\item Local UnsolvedMath registry, permanent ID \(30006200\), source code \texttt{catalogue entry}, title ``Optimization Workshop Contents Fragment,'' in the collection ``Oberwolfach Reports Open Problems.'' The supplied record gives proposed year 2025 and no identified problem proposer.
\item Supplied report citation: \emph{MFO-RIMS Tandem Workshop: Optimization, Theoretical Computer Science and Algebraic Geometry: Convexity and Beyond} (2025), \emph{Oberwolfach Reports}, pp.\ 405--450, DOI 10.4171/owr/2025/9. Supplied source URL and literature-source link labeled ``OWR report'': \url{https://doi.org/10.4171/owr/2025/9}. The supplied background describes Oberwolfach Reports as documenting mathematical workshops and their open-problem discussions. This report-level citation does not identify a recoverable conjecture in the supplied extraction.
\item Fragment attribution, p.\ 416: Luis Felipe Vargas, joint with Alex Bortolotti and Monaldo Mastrolilli, ``On the degree automatability of SoS proofs.'' The preceding isolated ``Conjecture'' entry also gives p.\ 416, without a recoverable assertion or explicit attribution.
\item Fragment attribution, p.\ 419: Dmitriy (Tim) Kunisky, joint with Maximilian Jerdee, Cristopher Moore, and Alexander S. Wein, ``Tensor networks, tensor cumulants, tensor PCA.''
\item Fragment attribution, p.\ 421: Petter Brändén, joint with Leonardo Saud Maia Leite, ``Totally nonnegative matrices, chain enumeration and zeros of polynomials.''
\item Fragment attribution, p.\ 423: Sebastian Debus, joint with Jose Acevedo, Grigoriy Blekherman, and Cordian Riener, ``Symmetric nonnegative functions and the Vandermonde map.''
\item Fragment attribution, p.\ 424: Sandra Kiefer, joint with Brendan D. McKay and Devini de Mel, ``Hard Graphs for Colour Refinement and Beyond.''
\item Fragment attribution, p.\ 425: Grigoriy Blekherman, ``Graph Homomorphisms and Polynomials.''
\item Fragment attribution, p.\ 426: Yoshio Okamoto, ``Combinatorial Shorest Paths on Combinatorial Polytopes'' (title spelling as supplied).
\item Fragment attribution, p.\ 428: Shin-ichi Tanigawa, ``Identifying Generic Points from Non-generic Measurements.''
\item Fragment attribution, p.\ 429: Keiya Sakabe, joint with Hiroshi Hirai, ``Gradient descent for unbounded convex functions on Hadamard manifolds and its applications to scaling problems.'' These contributor attributions concern contents entries and do not establish authorship of an identifiable conjecture.
\item Supplied literature assessment, checked on \texttt{2026-08-22}: the extraction is a table-of-contents fragment with no identifiable mathematical proposition, so its truth cannot be tested and it is retained as ``open.'' This is supplied dated curation derived from the supplied OWR material, not independently verified here. The supplied background expressly calls for independent verification of time-sensitive status and cited claims.
\end{itemize}
\end{document}
